\bookchapter{Capitolo 34: Guida approfondita alla connettività auto-telefono — la fusione tra telefono e automobile}

\chapterintro{L'auto sta evolvendo da “quattro ruote più un divano” a “quattro ruote più uno smartphone”. Che tu guidi un'auto a combustione o a nuova energia, il rapporto tra telefono e infotainment determina l'esperienza di guida. Questo capitolo ti aiuta a chiarire le tre grandi vie della connettività auto-telefono, le caratteristiche dei sistemi infotainment dei vari schieramenti e quei piccoli trucchi che rendono la guida più scorrevole.}

\booksection{Le tre grandi vie della connettività auto-telefono}

Attualmente ci sono tre approcci principali per “connettere” telefono e auto; capito quelli, si capisce l'intero ecosistema:

\begin{itemize}
  \item \textbf{Proiezione del telefono sull'infotainment}: l'infotainment funge solo da "schermo" e "altoparlante", calcolo e dati cartografici provengono dal telefono. Soluzioni rappresentative: Apple CarPlay, Google Android Auto, Huawei HiCar, Baidu CarLife. Il vantaggio è che l'ecosistema del telefono viene trasferito direttamente in auto senza reinstallare le app; lo svantaggio è che le capacità hardware dell'infotainment restano inutilizzate.
  \item \textbf{Integrazione profonda telefono-infotainment}: l'infotainment è un sistema indipendente, ma condivide account, app e dati col telefono, e può persino "passare il testimone" dei compiti. Soluzioni rappresentative: abitacolo HarmonyOS (Aito/Luxeed), BYD DiLink, NIO OS di NIO, sistema infotainment di Li Auto. Il vantaggio è l'esperienza integrata; lo svantaggio è il legame allo stesso ecosistema di marca.
  \item \textbf{CarPlay 2.0 versione profondamente integrata}: la nuova generazione di CarPlay presentata da Apple nel 2025 non è più solo "proiezione", ma prende in carico il cruscotto, il controllo del climatizzatore, la regolazione del sedile e altre funzioni chiave dell'infotainment. Ferrari, Aston Martin, Porsche e altri marchi lo montano per primi, segno che il sistema infotainment entra in una nuova fase "guidata dai produttori di smartphone".
\end{itemize}

\booksection{Confronto CarPlay / Android Auto / HiCar}

I tre principali schemi di proiezione hanno punti di forza diversi; quando scegli auto o telefono puoi consultare la tabella seguente:

\begin{center}
\begin{tabularx}{\linewidth}{|X|X|X|X|X|}
\hline
\textbf{Soluzione} & \textbf{Sistema telefonico supportato} & \textbf{Modelli auto supportati} & \textbf{Ricchezza di funzioni} & \textbf{Usabilità in Cina} \tabularnewline
\hline
Apple CarPlay & solo iPhone & molto ampia, supportata dalla maggior parte dei modelli & alta, mappe/musica/telefono/messaggi ok & buona, nativamente disponibile \tabularnewline
Android Auto & Android (telefoni cinesi richiedono framework servizi) & abbastanza ampia, soprattutto joint venture/import & alta, paragonabile a CarPlay & media, richiede attivazione via VPN \tabularnewline
Huawei HiCar & principalmente Huawei, parziale supporto altri Android & auto cinesi e alcune joint venture & abbastanza alta, può interfacciarsi con abitacolo HarmonyOS & molto buona, ecosistema locale \tabularnewline
Baidu CarLife & sia iOS che Android & abbastanza ampia, comune su modelli vecchi post-installati & media, aggiornamenti lenti & molto buona, ecosistema locale \tabularnewline
\hline
\end{tabularx}
\end{center}

\booksection{Abitacolo HarmonyOS: passaggio senza soluzione di continuità tra telefono e auto}

Se CarPlay è “il telefono proiettato sull'infotainment”, l'abitacolo HarmonyOS è “telefono e infotainment sono una cosa sola”. La collaborazione tra telefoni Huawei e infotainment di Aito/Luxeed rappresenta il livello più alto attuale di connettività auto-telefono in Cina:

\begin{itemize}
  \item \textbf{Super Desktop}: l'infotainment richiama direttamente le app del telefono, senza reinstallarle. La musica in ascolto o il percorso di navigazione sul telefono si trasferisce automaticamente allo schermo dell'auto quando sali, e torna al telefono quando scendi, senza interruzioni.
  \item \textbf{Passaggio di applicazione}: la navigazione impostata sul telefono viene rilevata dall'infotainment alla chiusura della portiera; arrivato al parcheggio, dopo essere sceso, il telefono continua la navigazione pedonale fino alla destinazione finale. L'"ultimo chilometro" non si interrompe più.
  \item \textbf{Account condiviso telefono-auto}: un solo login dell'account Huawei e l'infotainment sincronizza automaticamente l'abbonamento musicale, i preferiti delle mappe, il calendario, i contatti. Infotainment e telefono sono terminali diversi della stessa identità, non due sistemi separati.
\end{itemize}

\booksection{I sistemi infotainment delle nuove forze dell'auto}

L'infotainment delle auto a nuova energia cinesi è passato da “navigazione + musica” a terminale intelligente mobile, e ogni marca ha le sue caratteristiche distintive:

\begin{itemize}
  \item \textbf{NIO OS di NIO}: l'assistente vocale NOMI è un punto di riferimento del settore, riconosce comandi multipli e dialoghi contestuali continui, ed esprime anche emozioni con faccine. NIO Phone trasforma il telefono in "chiave dell'auto + secondo schermo", legato in profondità all'ecosistema NIO.
  \item \textbf{Infotainment Li Auto}: il punto forte è l'interazione multi-schermo: lo schermo del passeggero può riprodurre serie TV in modo indipendente, quello per i bambini può mostrare cartoni animati. Il "Task Master" permette di personalizzare gli scenari, ad esempio "metti in pausa la musica aprendo la portiera, riprendi chiudendola".
  \item \textbf{Infotainment Xiaomi HyperOS}: interagisce in profondità con telefono Xiaomi, orologio ed elettrodomestici. La sveglia impostata sul telefono viene ricordata dall'infotainment; quando navighi verso casa, il climatizzatore di casa si accende in anticipo. HyperOS unifica l'"ecosistema completo persona-auto-casa".
  \item \textbf{BYD DiLink}: caratteristico schermo rotante, verticale per i brevi video, orizzontale per le mappe, si adatta all'orientamento dell'app. Include un mercato app con una ricchezza di ecosistema leader tra i costruttori tradizionali.
\end{itemize}

\booksection{Il telefono come infotainment — si può usare?}

Se il tuo infotainment è troppo vecchio o troppo brutto, puoi usare direttamente il telefono “al suo posto”? La risposta è: sì, ma dipende dallo scenario.

\begin{itemize}
  \item \textbf{Motorola Ready For / Modalità desktop}: alcuni telefoni Motorola, Samsung, Huawei supportano la "modalità desktop", collegati a un monitor diventano un computer elementare. In teoria possono funzionare anche con lo schermo dell'auto, ma la compatibilità reale è discontinua e mancano interazioni auto-specifiche (come i tasti sul volante).
  \item \textbf{Supporto telefono vs infotainment originale}: se l'infotainment originale scatta o le mappe non si aggiornano, un solido supporto per telefono + navigazione sul telefono è una scelta pragmatica. Ma attenzione: (1) lo schermo del telefono è piccolo, in guida non si vede bene; (2) in estate il sole caldo può surriscaldare il telefono con throttling o spegnimento automatico; (3) la navigazione prolungata consuma molta batteria, serve un caricabatteria da auto.
  \item \textbf{Consiglio}: per tragitti brevi o uso occasionale, il supporto telefono basta; per il pendolarismo quotidiano e i lunghi viaggi, conviene aggiornare a un infotainment che supporti CarPlay/HiCar o installarne uno after-market, per esperienza e sicurezza migliori.
\end{itemize}

\begin{tcbwarn}{Promemoria di sicurezza}
Qualsiasi soluzione si usi, usare il telefono durante la guida è un'infrazione al codice della strada. Il supporto telefono serve solo per “guardare la navigazione”, non per “operare il telefono”. Impostare la destinazione, cambiare brano e simili va fatto da fermo, oppure col controllo vocale.

\end{tcbwarn}

\booksection{Consigli pratici}

\begin{itemize}
  \item \textbf{CarPlay wireless instabile}: (1) spegni il Bluetooth dell'infotainment e riparetalo; (2) elimina i dispositivi abbinati sull'infotainment e ricollegali; (3) controlla se il telefono è connesso a più dispositivi Bluetooth contemporaneamente causando conflitti; (4) aggiorna il firmware dell'infotainment e il sistema del telefono all'ultima versione; (5) in caso contrario, passa a CarPlay cablato, più stabile.
  \item \textbf{Traffico dati dell'infotainment esaurito}: (1) apri un hotspot col telefono e collega l'infotainment a quello; (2) nelle impostazioni dell'infotainment sposta su Wi-Fi le app ad alto consumo (come video e musica in streaming); (3) pulisci periodicamente la cache delle mappe offline dell'infotainment.
  \item \textbf{Attivazione automatica della modalità di guida}: (1) la "modalità Concentrazione" di iOS può essere impostata per attivarsi automaticamente "alla connessione al Bluetooth dell'auto", silenziando le notifiche non urgenti; (2) la "modalità di guida" di Android o app di terzi (come Tasker) possono silenziare e rispondere automaticamente ai messaggi alla connessione Bluetooth; (3) i telefoni Huawei possono impostare il "pagamento intelligente" per passare automaticamente alla riproduzione musicale sull'auto quando si sale.
\end{itemize}

\begin{tcbinfo}{Piccoli consigli nella scelta dell'auto}
Se tieni all'esperienza dell'infotainment, in concessionaria prova con mano: (1) porta il tuo telefono e prova la stabilità della connessione CarPlay/HiCar; (2) prova la velocità di reazione e il tasso di riconoscimento comandi dell'assistente vocale; (3) controlla se il mercato app ha le applicazioni che usi di solito; (4) informati sulla frequenza degli aggiornamenti OTA futuri — un buon infotainment si costruisce con l'iterazione software.

\end{tcbinfo}

\begin{tcbpractice}{Fai: ottimizza la tua esperienza di connettività auto-telefono}
\begin{enumerate}
  \item Verifica quale schema di interconnessione supportano il tuo telefono e la tua auto (CarPlay / HiCar / Android Auto / CarLife).
  \item Pulisci una volta l'elenco delle associazioni Bluetooth dell'infotainment, tenendo solo i telefoni abituali, per ridurre i conflitti di connessione.
  \item Imposta una "modalità Concentrazione di guida" che si attivi alla connessione del Bluetooth dell'auto.
  \item Scarica l'app del marchio dell'auto (come HarmonyOS Intelligent Mobility, Li Auto, NIO App) e verifica se puoi controllare da remoto il climatizzatore e vedere lo stato dell'auto.
\end{enumerate}

\end{tcbpractice}

\begin{tcbsummary}{Riepilogo del capitolo}
\begin{itemize}
  \item Tre vie della connettività auto-telefono: proiezione (CarPlay ecc.), integrazione profonda (abitacolo HarmonyOS ecc.), CarPlay 2.0 profondamente integrato
  \item CarPlay ha l'ecosistema più ampio; HiCar offre la migliore esperienza in Cina; Android Auto in Cina richiede interventi; CarLife è adatto a post-installazione su auto vecchie
  \item Il Super Desktop, il passaggio di applicazione e l'account condiviso dell'abitacolo HarmonyOS rappresentano la direzione dell'integrazione profonda
  \item NOMI di NIO, il multi-schermo di Li Auto, l'ecosistema persona-auto-casa di Xiaomi e lo schermo rotante di BYD hanno ciascuno le proprie caratteristiche
  \item Il telefono come infotainment è adatto solo a scenari temporanei; per il lungo periodo conviene un upgrade o un after-market
  \item Consigli pratici: riconnessione CarPlay, hotspot del telefono per il traffico, modalità di guida attivata via Bluetooth
\end{itemize}

\end{tcbsummary}
